% ============================================================
% corex-header-footer.tex
% Header & footer terpisah dari preamble supaya bisa diubah
% independen (misal ganti teks footer tanpa sentuh styling lain).
% Butuh \CorexDocTitle dan \CorexClientName didefinisikan SEBELUM
% \input file ini (lihat contoh di examples/).
% ============================================================

\pagestyle{fancy}
\fancyhf{}
\renewcommand{\headrulewidth}{0.6pt}
\renewcommand{\footrulewidth}{0.4pt}
\renewcommand{\headrule}{\hbox to\headwidth{\color{corexsignal}\leaders\hrule height \headrulewidth\hfill}}

% ---------- Header ----------
\fancyhead[L]{\small\bfseries\color{corexnavy} COREX}
\fancyhead[R]{\small\color{corexgray} \CorexDocTitle}

% ---------- Footer ----------
\fancyfoot[L]{\footnotesize\color{corexgray} Confidential --- Disiapkan untuk \CorexClientName}
\fancyfoot[R]{\footnotesize\color{corexgray} Halaman \thepage\ dari \pageref{LastPage}}

% Halaman pertama (title page) tidak pakai header/footer standar
\fancypagestyle{corextitle}{
  \fancyhf{}
  \renewcommand{\headrulewidth}{0pt}
  \fancyfoot[C]{\footnotesize\color{corexgray} Confidential --- Disiapkan untuk \CorexClientName}
}
